\documentclass[11pt]{article}
\usepackage{amsmath,amssymb,amsthm,booktabs}
\usepackage[margin=2.4cm]{geometry}
\usepackage[hyphens]{url}\sloppy
\newtheorem{theorem}{Theorem}\newtheorem{lemma}[theorem]{Lemma}\newtheorem{proposition}[theorem]{Proposition}
\newtheorem{corollary}[theorem]{Corollary}
\theoremstyle{definition}\newtheorem{definition}[theorem]{Definition}\newtheorem{remark}[theorem]{Remark}
\newcommand{\Z}{\mathbb Z}\newcommand{\Oth}{\operatorname{Oth}}
\title{Q597 (Share$_2$): the second handshake is forced by the tiling\\
\normalsize PROVE 2026-11-26 c2. \textbf{(Share$_2$) is proved}, with one hypothesis on $E$ (H1: $J_E$ takes a step after its contact), from Theorem~share, Proposition~orient, Lemma~word, the finality of created tiles (Lemma~\ref{lem:fin}, proved here from MU, REM and CRE) and one finite read of the configuration table (Lemma~\ref{lem:loc2}). The monotone-cycle method of the share tex was \emph{not} needed. The killing constraint is the plainest one available: $M_0$ is a tiling. All the tiles removed at the second step are $M_0$ tiles. The contact hexagons kill each role's own a-type at step~1. Of the four remaining pairings, the two with equal letters (a with a, b with b) put two overlapping squares on $M_0$. The two surviving pairings share their two outside tiles. So the shapes of $E_0,E_1,L_0,L_1$ form a joint classification with 8 classes, proved. It coincides with the census: 779/779 records at size~5, and 718/718 on fresh size-6 chunks 11--12. Lesion: without tiling compatibility, 12 combinations survive and (Share$_2$) fails on 4. \textbf{The same mechanism proves the corridor} (Theorem~\ref{thm:corr}, the $M_0$ half of (Merge)). At a shared class-0 rhombus, $B$ and $C$ remove the same tiles for each type, and stopping is target-blind. Both journeys remove only $M_0$ tiles, so Lemma~ovl forces equal types. The journeys then take identical a-steps until a common $(0,b)$, where the two targets rotate differently. That is the only place $B\neq C$ is used. In case~b, (F3) of the WSlast tex (proved) gives (H1), the exit, and $k=m_L-3$ (Proposition~\ref{prop:shadow}). So both theorems hold there with no extra hypothesis, and the corridor has two independent proofs. Still computed: no shared step after the exit (this clause fails on all 249 $B/B$ records that pass the exit), and (H1) in case~a. (Share$_2$) itself does \emph{not} separate $B/C$ from $B/B$: it holds on 440/440 $B/B$ pairs. DREAM's zone falsifier came back \textbf{negative}.}
\author{Clio}\date{26 November 2026 (second cycle)}
\begin{document}\maketitle

\section*{Provenance}
\begin{itemize}
\item Code: \texttt{proofs/code-1126c2-Q597/} (new; no script shadows a stdlib module). \texttt{PREDICTION.md} was written before any of the following was run. \texttt{step0s2.py} (Step~0 predicates, translation classes and the zone test, reading the record pickles of \texttt{code-1126-Q597/step0.py}); \texttt{local2.py} (the local classification, with \texttt{LESION}); \texttt{why2.py} (the killing constraint one role at a time, the census match, and the local $B/B$ control); \texttt{step0BB.py} (a copy of \texttt{step0.py} that records non-commuting $B/B$ pairs) and \texttt{bbshare2.py} (the predicate on them). Each writes \texttt{X.out}.
\item Engine and census: \texttt{code-1126-Q597/step0.py} through the \texttt{code-1125c2-Q597/bsall.py} preamble, unchanged. Size~5 is the full census. Size~6 uses chunks 11 and 12 of 24. They are fresh: c1 used 7--10 and \texttt{mos6} used 3--6.
\item Configuration geometry: \texttt{code-1122-Q597/geom.py} (the AW table \texttt{words.configs}; overlap by the separating-axis test on exact integer points, tolerance $10^{-9}$), unchanged.
\item Upstream, read first-hand this session: \texttt{2026-11-26-Q597-G3prime.tex} (all), \texttt{2026-11-23-Q597-share.tex} (all), \texttt{2026-11-22-Q597-first-contact.tex} (Lemma~U, Lemma~dinv), \texttt{2026-11-20-Q597-no-repeat.tex} (the engine rule, MU, REM, CRE, Lemma~class), \texttt{2026-11-21-Q597-G1a-A1.tex} (Lemma~nr), \texttt{2026-11-25c2-Q597-BS-assembly.tex} (Theorem~bsall), \texttt{2026-11-22c2-Q597-role.tex} (Proposition~orient). No external source. No browsing.
\end{itemize}

\section{Setting}\label{sec:set}
As in \S1 of the G3$'$ tex. $(x,y)$ is a diverging $B/C$ pair at a census state $M_0$, a Purbhoo mosaic. $L$ is the role whose contact step is the class-1 a-step $(1,a_{T_L})$ of its own target (Proposition~orient, unconditional since Theorem~share), and $E$ is the other role. Write $T_E=T_L'$. $p_E,p_L$ are the contact (= divergence, Theorem~fc) steps, and $E_i=J_E[p_E+i]$, $L_i=J_L[p_L+i]$. A step $s$ removes the four tiles $H(s)$ (its rhombus and $G(s)$, one square and two triangles) and creates the four tiles $C(s)$, one of which is the next rhombus. Put
\[H_E=H(E_0),\quad H_L=H(L_0),\qquad N=\bigl(H_E\cup H_L\bigr)\setminus\{\text{rhombi}\},\qquad \Oth(E_1)=G(E_1)\setminus N,\quad \Oth(L_1)=G(L_1)\setminus N.\]
\emph{Classes} are the absolute ones of the no-repeat tex, Lemma~class: a rhombus is $s_0\times s_1$ with both edges parallel to one vector $u\in\{(0,1),(1,0),(1,-1)\}$. The class $u$ is allowed for $T$ iff $\eta_T(u)\neq0$. For $B$ the allowed classes are $(0,1)$ and $(1,-1)$, and for $C$ they are $(0,1)$ and $(1,0)$. Class~0 is $(0,1)$, the one class allowed for both targets. Class~1 is the other allowed class of the target in question.
\begin{remark}[the census's relative class is this one]
The census scripts call class~0 ``the class of the journey's first rhombus''. Every nest-$A$ corner starts a journey toward $B$ and one toward $C$, since pairs are ordered both ways. So under termination its class is allowed for both targets, and it is $(0,1)$. The census confirms the two conventions agree. All 1497 records of \S\ref{sec:step0} fall into the eight classes of Lemma~\ref{lem:loc2}, which are computed in the absolute convention.
\end{remark}
\paragraph{Standing hypotheses} (S), as in the 11-25 c2 tex: termination, Purbhoo mosaic, engine rule (S1)--(S4). Also:
\begin{itemize}
\item[(H1)] $J_E$ takes a step at $p_E+1$.
\end{itemize}
For $L$ the corresponding fact is part of BS$(L)$ (Theorem~bsall, proved). $L$'s middle is $a\cdot C$ with $C$ nonempty, so $L_1$ exists. (H1) is computed, not proved: $|{\rm mid}_E|\ge3$ on every record (the (LEN) law, G3$'$ tex). It is the only place where a fact about $E$ enters beyond its contact.

\section{Step 0 (computed)}\label{sec:step0}
\paragraph{The predicate} (\texttt{PREDICTION.md}, written before running).
\begin{itemize}
\item[(S2.i)] $G(E_1)$ contains no tile of $C(E_0)$, exactly one tile of $H_L$ and none of $H_E$. The same holds for $L$ with $E$ and $L$ exchanged.
\item[(S2.ii)] $\Oth(E_1)=\Oth(L_1)$, and it has two tiles.
\item[(S2.iii)] $|C(E_1)\cap C(L_1)|=2$ in case~b and $0$ in case~a.
\end{itemize}
\begin{center}\small\begin{tabular}{lrrrr}\toprule
& size 5 a & size 5 b & size 6 (chunks 11--12) a & size 6 b\\\midrule
records & 258 & 521 & 277 & 441\\
(S2.i) both roles & 258 & 521 & 277 & 441\\
(S2.ii), $|\Oth|=2$ & 258 & 521 & 277 & 441\\
$\Oth$ kinds & \{sq, tri\} & \{sq, tri\} & \{sq, tri\} & \{sq, tri\}\\
the $H_L$ tile of $E_1$, the $H_E$ tile of $L_1$ & tri, tri & tri, tri & tri, tri & tri, tri\\
$|C(E_1)\cap C(L_1)|$ & 0 & 2 (rh, tri) & 0 & 2 (rh, tri)\\
joint classes of $(E_0,E_1,L_0,L_1)$ & 4 & 4 & 4 & 4\\\bottomrule
\end{tabular}\end{center}
The predicted kinds (one square and one triangle outside, a triangle of the other hexagon) were right. So was the class count: 4 in case~b (the (W$_F$) contact types) and 4 in case~a. The case-a classes are the prefixes of the (W$_a$) classes, and the two size-6 ``new'' (W$_a$) classes of c1 already have their prefix here at size~5 (1 and 8 records). In case~b the common created tiles are a rhombus and a triangle. So \emph{$E_2$ and $L_2$ start from the same rhombus}. That is where the corridor begins.

\paragraph{(b) DREAM's zone falsifier: negative.} $\alpha$ is a lattice unit vector on all 256 corridor records at size~5. The prediction was that every tile removed by a corridor step of $E^L$, of $L^E$ or of $J_E$ has an edge parallel to $\alpha$, and that these tiles are joined along $\alpha$-edges. Both parts fail. On \emph{no} record do all corridor tiles have an $\alpha$-edge (typically 2 or 3 of the 4 tiles of a step do). The $\alpha$-edge tiles are connected along shared $\alpha$-edges in 97 of the 256 records and disconnected in 159. The counts are the same as the 97/159 split of the G3$'$ corridor lesion, where 97 records have a corridor that contributes $0$ to each $\Delta$. I did \emph{not} check that they are the same records. So the corridor is \emph{not} one de Bruijn zone along $\alpha$ in the sense tested, and the train reading of (Merge$'$) loses its cheapest support. The coincidence of counts is recorded, not explained.

\section{Created tiles are final}\label{sec:fin}
\begin{lemma}[finality]\label{lem:fin}
In a terminating journey $J$ toward $T$ in a Purbhoo mosaic, no non-rhombus tile created by a step is removed by a later step. Consequently every non-rhombus tile that $J$ removes is a tile of the initial board, and $J$ removes it exactly once.
\end{lemma}
\begin{proof}
Let $t$ be created at step $k$ and removed at step $j>k$. By CRE (no-repeat tex, proved), $\psi_{T,i}(t)\ge(\mu_{k+1})_i+1$ for some $i$. By REM, $\psi_T(t)\le\mu_j$ componentwise. By MU, $\mu$ is componentwise non-increasing, so $\mu_j\le\mu_{k+1}$. Hence $\psi_{T,i}(t)\le(\mu_j)_i\le(\mu_{k+1})_i<\psi_{T,i}(t)$, a contradiction.
A tile on the board before step $j$ is either an initial tile not yet removed or a tile created by an earlier step. A removed non-rhombus tile is not the latter, so it is an initial tile. If it were removed twice, it would have been re-created in between, which is the case just excluded (and also Lemma~nr of the A1 tex).
\end{proof}
\noindent This is the ``no revisit'' used informally in the share tex (``$\mathrm{rem}_x(u)$ is unique, by no-revisit and Lemma~nr''). Here it is stated and proved.

\begin{corollary}[the step-1 tiles are $M_0$ tiles]\label{cor:m0}
Assume (S) and (H1). Then $N\cup G(E_1)\cup G(L_1)\subseteq M_0$. Moreover $G(E_1)\cap(C(E_0)\cup H_E)=\emptyset$ and $G(L_1)\cap(C(L_0)\cup H_L)=\emptyset$. Hence, since $M_0$ is a tiling, any two tiles of $N\cup G(E_1)\cup G(L_1)$ are equal or have disjoint interiors.
\end{corollary}
\begin{proof}
Lemma~\ref{lem:fin} for $J_E$ and $J_L$ on $M_0$. The non-rhombus tiles of $H_E$ and $G(E_1)$ are removed by $J_E$, so they are $M_0$ tiles, each removed once. Hence $G(E_1)\cap H_E=\emptyset$, and $G(E_1)\cap C(E_0)=\emptyset$ by finality. The same argument applies to $L$, using $L_1$ from BS$(L)$.
\end{proof}

\section{The local classification (proved; finite part exact)}\label{sec:loc}
\begin{lemma}[the joint shapes at steps 0 and 1]\label{lem:loc2}
Let $T_L\in\{B,C\}$ and $T_E=T_L'$. Consider quadruples $(e_0,\ell_0,e_1,\ell_1)$ of good configurations (Lemma~word) such that:
\begin{enumerate}
\item $\ell_0$ is the $(1,a_{T_L})$ configuration of $T_L$ at some rhombus $R_L$. $e_0$ is a good $T_E$ configuration at some rhombus $R_E$, other than $(1,a_{T_E})$. The two have the same square $u$.
\item $e_1$ is a good $T_E$ configuration at the rhombus created by $e_0$, and $\ell_1$ is a good $T_L$ configuration at the rhombus created by $\ell_0$.
\item (tiling) any two tiles of $N\cup G(e_1)\cup G(\ell_1)$ are equal or interior-disjoint, where $N$ is the set of non-rhombus tiles of $H(e_0)\cup H(\ell_0)$.
\end{enumerate}
Up to a common translation there are exactly 8 such quadruples. On each of them:
\begin{itemize}
\item[(a)] $e_1$ is $(0,a_{T_L})$ or $(0,b)$, and $\ell_1$ is $(1,b)$ or $(1,a_{T_E})$. \textbf{Exactly one of $e_1,\ell_1$ is a b-step}: $(e_1,\ell_1)\in\{((0,a_{T_L}),(1,b)),\ ((0,b),(1,a_{T_E}))\}$, and for each contact pair both occur.
\item[(b)] $G(e_1)$ contains exactly one tile of $H(\ell_0)$, a triangle, and no tile of $H(e_0)$. Symmetrically for $\ell_1$.
\item[(c)] $G(e_1)\setminus N=G(\ell_1)\setminus N$, consisting of one square and one triangle.
\item[(d)] $|C(e_1)\cap C(\ell_1)|$ is $2$ (a rhombus and a triangle) when $\ell_1=(1,b)$, and $0$ when $\ell_1=(1,a_{T_E})$.
\end{itemize}
\end{lemma}
\begin{proof}[Proof (finite, exhaustive, exact)]
Translate so that $u$ is fixed. Then condition~1 fixes $\ell_0$, and $e_0$ ranges over the five allowed $T_E$ configurations whose square has the shape of $\ell_0$'s square (one translate each). Two of the five qualify for each $T_L$ (\texttt{local2.out}, \texttt{why2.out}): $(1,b)$ and $(0,a_{T_E})$. Condition~2 leaves three types for each of $e_1,\ell_1$, so there are $2\times2\times3\times3=36$ quadruples. The square-sharing contact hexagons satisfy condition~3 on their own in all four contact pairs. Then:
\begin{itemize}
\item \emph{One role at a time.} For each contact pair, $e_1=(0,a_{T_E})$ violates condition~3: its square overlaps two triangles of $N$. Likewise $\ell_1=(1,a_{T_L})$: its square overlaps one or two triangles of $N$. The other two types of each role pass (\texttt{why2.out}, lines \texttt{E1 type}, \texttt{L1 type}). So each role's own a-type is excluded at step~1.
\item \emph{The pair.} Of the four remaining pairings, $(e_1,\ell_1)=(a,a)$ and $(b,b)$ violate condition~3: the two squares overlap (and a square overlaps a triangle in each direction). The pairings $(a,b)$ and $(b,a)$ pass (\texttt{PAIR} lines).
\end{itemize}
That leaves $4\times2=8$ quadruples. (a)--(d) are read off each of them (\texttt{local2.out}, \texttt{SURV} lines: Oth equal, $|\Oth|=2$, $|G(e_1)\cap H(\ell_0)|=|G(\ell_1)\cap H(e_0)|=1$, no meeting with created tiles, and $|C\cap C|=2,0,2,0,\dots$). The classes of $e_1,\ell_1$ in (a) follow from Lemma~class: an a-step keeps the class and a b-step changes it. So $e_1$ sits at class~0 in both contact types, and $\ell_1$ at class~1.
\end{proof}
\noindent\emph{Instrument controls.} (i) The overlap test fires on the excluded types and stays silent on the surviving ones. Equal tiles are treated as compatible, and that case occurs: the shared square is in $N$ from both hexagons, and the $H_L$-triangle of $e_1$ is equal to a tile of $N$. (ii) \texttt{local2.py} also evaluates the board-level constraints: the boards of $J_E$ before and after $E_0$, those of $J_L$, and Lemma~U at both rhombi and both contacts. With these \emph{added}, the same 8 survive. With \emph{only} condition~3 kept (\texttt{local2\_M0only.out}), the same 8 survive. So condition~3 is the whole killing constraint.

\section{The theorem}\label{sec:thm}
\begin{theorem}[(Share$_2$): the second handshake]\label{thm:share2}
Assume (S) and (H1), and let $(x,y)$ be a diverging $B/C$ pair with roles $L,E$ as in \S\ref{sec:set}. Then:
\begin{enumerate}
\item $\Oth(E_1)=\Oth(L_1)$. It consists of one square and one triangle, both $M_0$ tiles outside both contact hexagons.
\item $G(E_1)$ contains exactly one tile of $H_L$, a triangle, and $G(L_1)$ contains exactly one tile of $H_E$, a triangle. $E_1$ removes no tile of $H_E\cup C(E_0)$ other than its rhombus, and $L_1$ none of $H_L\cup C(L_0)$.
\item Exactly one of $E_1,L_1$ is a b-step. In case~b, $L_1=(1,b)$ and $E_1=(0,a_{T_L})$, and $C(E_1)\cap C(L_1)$ consists of a rhombus $R^*$ of class~0 and a triangle. So $E_2$ and $L_2$, whenever they exist, are taken from the same rhombus $R^*$. In case~a, $L_1=(1,a_{T_E})$, $E_1=(0,b)$, and $C(E_1)\cap C(L_1)=\emptyset$.
\end{enumerate}
\end{theorem}
\begin{proof}
By Lemma~word each of $E_0,E_1,L_0,L_1$ is a translate of a good configuration of its own target at its rhombus. By the definition of a journey, $E_1$ is taken at the rhombus created by $E_0$, and likewise for $L$. By Proposition~orient, $L_0=(1,a_{T_L})$ and $E_0\ne(1,a_{T_E})$. By Theorem~share, $H_E$ and $H_L$ have the same square. By Corollary~\ref{cor:m0}, condition~3 of Lemma~\ref{lem:loc2} holds. So $(E_0,L_0,E_1,L_1)$ is a translate of one of the eight quadruples of Lemma~\ref{lem:loc2}, and (a)--(d) there give items 1, 2 and the first sentence of~3 (the translate of $N$ is the $N$ of the pair).
For the cases: by Theorem~bsall, $L$'s re-run middle is $(C-\alpha)\cdot(a+\gamma)$ with $C=(L_1,\dots)$, so the re-run's step at $p_L$ is $L_1-\alpha$. A translate has the same type and class. In case~b that step is $(1,b)$, so $L_1=(1,b)$. In case~a, $C$ is the single step $(1,a_{T'})=(1,a_{T_E})$ (Theorem~bsall). Lemma~\ref{lem:loc2}(a),(d) give the rest. $R^*$ is a rhombus of class allowed for both targets, hence of class~0. Directly: $E_1$ is an a-step at class~0, so it keeps class~0.
\end{proof}

\begin{corollary}[the front of (W$_F$) and (W$_a$), on $M_0$]\label{cor:front}
The $M_0$-half of the front windows, $(E_0,E_1,L_0,L_1)$, is proved to take one of 8 shapes up to translation: 4 in case~b and 4 in case~a. This proves the journey columns $E_0,E_1,L_0,L_1$ of (W$_F$), and the corresponding prefix of (W$_a$). The re-run columns $E^L_0,E^L_1,L^E_0$ are not covered by Theorem~\ref{thm:share2}. For $L^E_0$, BS$(L)$ supplies them. $E^L_0$ and $E^L_1$ remain open.
\end{corollary}

\section{Controls}\label{sec:ctl}
\begin{itemize}
\item \textbf{Out of sample.} The local classification is exhaustive, not a sample of a census. It needs no extension at size~7: no input depends on the size. The census is corroboration only. Size~5 full: 779/779 records lie in the 8 classes, and every class is hit. Size-6 chunks 11--12: 718/718, every class is hit (\texttt{why2.out}, \texttt{CENSUS} lines; exact tile-set equality after translating the shared square to its base point, created tiles included).
\item \textbf{Lesion of the hypothesis used.} Drop condition~3 (all tiling compatibility, \texttt{LESION=M0,XE0,XL0,XE1,XL1}): 12 quadruples survive the remaining (Lemma~U) constraints, and (Share$_2$) fails on 4 of them. Drop everything (\texttt{local2\_lesion\_all.out}): 36 survive and (Share$_2$) fails on 28. The control (no lesion) printed first: 8 survivors, (Share$_2$) on all 8.
\item \textbf{$B/B$, local.} The same lemma with $T_E=T_L$ (\texttt{why2\_BB.out}). Only 2 contact pairs share a square (the $(1,b)$ contact of $E$ needs the other target). On both, the same mechanism leaves $(a,b)$ and $(b,a)$, and (Share$_2$) holds. \textbf{So (Share$_2$) is not what separates commuting from non-commuting pairs.} It is a geometric fact about any pair whose $L$ role sits at $(1,a_T)$. The brief asked that it fail on $B/B$ ``or it is not doing work''. Locally it does not fail. Its work is to fix the shapes, not to make the boards commute. The $B/C$ vs.\ $B/B$ difference lies above it: the G3$'$ tex finds (LEN) and (Merge) failing on all 440 $B/B$ records.
\item \textbf{$B/B$, census} (\texttt{step0BB.py}, \texttt{bbshare2.out}; size~5, all 440 non-commuting $B/B$ pairs, each with exactly one $L$). (Share$_2$) \textbf{holds on 440/440}: (S2.i), (S2.ii), and $|C(E_1)\cap C(L_1)|=2$. There are 2 joint classes, with 241 and 199 records. \textbf{My prediction (c), that it would fail on more than half, was wrong.} This agrees with the local $B/B$ classification. (Share$_2$) is shape, not commutation.
\end{itemize}

\section{The corridor (Merge on $M_0$): proved, except one clause}\label{sec:corr}
The same mechanism carries past step~1. Three more inputs are needed.
\begin{lemma}[class~0 is target-blind; finite]\label{lem:c0}
Let $R$ be a rhombus of class~0. For each $t\in\{a0,a1,b\}$, $G^B_t(R)=G^C_t(R)$. For $t\in\{a0,a1\}$ the created sets also agree, so the $B$- and $C$-configurations of type $t$ at $R$ are the same step. For $t=b$ the created rhombi differ.
\end{lemma}
\begin{proof}
Read off the AW table at the canonical class-0 rhombus (exact points, \texttt{corr}-session check in \texttt{c0check.out}), and translate.
\end{proof}
\begin{lemma}[stopping is target-blind at a shared rhombus]\label{lem:pack}
Let $J,J'$ be journeys on $M_0$ from nest-$A$ corners, toward any targets. Suppose that before some step of each, neither has ended and both sit at the same rhombus $R$. Then $J$ ends there if and only if $J'$ does.
\end{lemma}
\begin{proof}
This is (S4) as the engine implements it (\texttt{code-1112c2-Q585/migrate.py}, \texttt{Board.journey}). Before each step the engine tests, for each nest $X$ other than the source nest $A$, whether $R$'s vertex set is a cell of $X$ at a position $ij$ not occupied in $X$ and such that the occupied set plus $ij$ is an order ideal. The test reads only $R$ and the position sets of the nests other than $A$. A journey changes a nest only when it ends. So before either has ended, both journeys see the nests of $M_0$, except $A$, which is excluded from the test. The inputs are equal, so the outcomes are equal.
\end{proof}

\begin{theorem}[the corridor]\label{thm:corr}
Assume (S) and (H1), and let the pair be in case~b. Put $R_2:=R^*$ (Theorem~\ref{thm:share2}(3)). There is $k\ge0$ such that:
\begin{enumerate}
\item for $2\le i\le k+1$, $E_i=L_i$ (identical steps). Each is a class-0 a-step taken by both journeys at the same rhombus $R_i$, and creates $R_{i+1}$;
\item at $R_{k+2}$ either both journeys end, or $E_{k+2}$ and $L_{k+2}$ are both $(0,b)$-steps with $H(E_{k+2})=H(L_{k+2})$ and different created rhombi;
\item if $E_j=L_{j'}$ for any steps $j,j'$ of the two journeys (indices relative to $p_E,p_L$, negative for prefixes), then either $j=j'\in\{2,\dots,k+1\}$, or $j,j'\ge k+3$.
\end{enumerate}
\end{theorem}
\begin{proof}
\emph{(1), (2).} Suppose that for some $i\ge2$ both journeys sit at a common class-0 rhombus $R_i$, with neither ended. This holds for $i=2$ by Theorem~\ref{thm:share2}(3). By Lemma~\ref{lem:pack}, both end there or both step. Suppose both step, $E$ with type $t$ and $L$ with type $t'$ (Lemma~word). The non-rhombus tiles they remove, $G^{T_E}_t(R_i)$ and $G^{T_L}_{t'}(R_i)=G^{T_E}_{t'}(R_i)$ (Lemma~\ref{lem:c0}), are all $M_0$ tiles (Lemma~\ref{lem:fin}). If $t\ne t'$, Lemma~ovl of the share tex says the square of $G^{T_E}_t(R_i)$ is not a tile of $G^{T_E}_{t'}(R_i)$ and its interior meets that of one of its tiles. Two distinct $M_0$ tiles would then overlap, which is impossible. So $t=t'$, and $H(E_i)=H(L_i)$. If $t$ is an a-type, Lemma~\ref{lem:c0} gives $E_i=L_i$, and the new rhombus $R_{i+1}$ is common. It has class~0, since an a-step keeps the class. Then repeat. If $t=b$, the created rhombi differ by Lemma~\ref{lem:c0}, since $T_E\neq T_L$. This is the only use of $B\neq C$ in the proof. That is (2) with $k=i-2$. The journeys terminate, so the repetition stops, and $k$ exists.
\emph{(3).} First, every step removes three non-rhombus tiles. If $E_j=L_{j'}$ these are removed by $J_L$, so they lie in $W_L$ (Lemma~\ref{lem:fin}). Hence $j<0$ would contradict the minimality of $p_E$: $p_E$ is the first step of $J_E$ whose hexagon meets $W_L$ (Theorem~fc). So $j\ge0$, and likewise $j'\ge0$. Second, by Lemma~\ref{lem:fin} a non-rhombus $M_0$ tile is removed by $J_L$ (or by $J_E$) at most once. Now go through the cases.
$j=0$: $E_0$ removes $u$, which $J_L$ removes only at $L_0$, so $j'=0$. But $E_0\ne L_0$: $L_0$'s rhombus has class~1 of $T_L$, which is not an allowed class of $T_E$.
$j=1$: $E_1$ removes a triangle of $H_L$ (Theorem~\ref{thm:share2}(2)), which $J_L$ removes only at $L_0$, so $j'=0$. But $L_0$ removes $u\in H_E$, and $E_1$ removes no tile of $H_E$.
$j'\in\{0,1\}$: the same two arguments with the roles exchanged.
$2\le j\le k+1$: $E_j=L_j$, so $L_{j'}$ and $L_j$ remove the same tiles, and $j'=j$. The case $2\le j'\le k+1$ is the same.
$j=k+2$ (or $j'=k+2$), if the journeys step there: $L_{j'}$ would remove the tiles $H(E_{k+2})=H(L_{k+2})$, so $j'=k+2$, but $E_{k+2}\ne L_{k+2}$ by (2).
What remains is $j,j'\ge k+3$.
\end{proof}

\paragraph{Census} (\texttt{corr5.out}, \texttt{corr6.out}; prediction appended to \texttt{PREDICTION.md} before the run). Every case-b record at size~5 (521) and size-6 chunks 11--12 (441) has the following properties. The shared a-steps are identical (406 and 488 steps). The exit is $(0,b)$ for both, with the same removed tiles and different created tiles. Both journeys stepping at the exit occurs in all records, and both ending in none. $k=m_L-3$. \textbf{No step after the exit is shared} (962/962). In case~a, $E_2$ and $L_2$ are never at the same rhombus (535/535).
\paragraph{$B/B$ control} (\texttt{corrBB5.out}). With $T_E=T_L$ the type-b configurations coincide too, so nothing makes the journeys part. On the 440 $B/B$ pairs the shared steps are identical throughout, b-steps included. Both journeys end together at a shared rhombus in 191 records (Lemma~\ref{lem:pack} in action). In the other 249, a shared b-step is identical and the journeys go on sharing steps, so the clause ``no shared step after the exit'' \textbf{fails 249/249}, and $k\neq m_L-3$ on all 249. So the part of (Merge) that separates $B/C$ from $B/B$ is the \emph{exit}: Lemma~\ref{lem:c0}'s different b-rotations, and what happens after them. Theorem~\ref{thm:corr}(1)--(3) are target-blind, except for the clause ``different created rhombi''.

\section{Case b without (H1): the shadow theorem already contains the corridor}\label{sec:shadow}
After Theorem~\ref{thm:corr} was written I re-read (F3) of the 11-25 WSlast tex, where it is a standing fact, \emph{proved for geometric records}. It comes from the sheet tex's Theorem~red with Proposition~hs and Corollary~geo (registry node \texttt{Q597-shadow-geometric}, proved). A geometric record is exactly the role $L$ in case~b, with $J_o=J_E$. In that notation (F3) reads as follows. $L$'s middle is $a\cdot C$ with $C=L[p_L+1..p_L+m]$ and $m\ge2$. $L_1=(1,b)$. For $2\le i\le m$, $L$'s rhombus before step $p_L+i$ equals $J_E$'s rhombus before step $c_0+i$, it has class~0, and $H(L_i)=H(J_E[c_0+i])$. $L_i$ is a class-0 a-step for $i<m$. $L_m$ is $(0,b)$ on the same hexagon as $J_E$'s $(0,b)$ step $e=c_0+m$. By Corollary~share, $c_0=c_0(L)=p_E$, so $J_E[c_0+i]=E_i$.
\begin{proposition}[case b, unconditional]\label{prop:shadow}
Assume (S), and let the pair be in case~b. Then:
\begin{enumerate}
\item (H1) holds. Hence Theorem~\ref{thm:share2} and Theorem~\ref{thm:corr} hold in case~b with no hypothesis beyond (S).
\item $E_i=L_i$ for $2\le i\le m-1$. $E_m$ and $L_m$ are both $(0,b)$ with $H(E_m)=H(L_m)$. So the corridor of Theorem~\ref{thm:corr} has $k=m-2$, and its exit is a step of both journeys, not a joint stop: this is (M1).
\item $m_L=m+1$, hence $k=m_L-3$: this is (M2).
\end{enumerate}
\end{proposition}
\begin{proof}
(1) $J_E$ has the step $c_0+m=p_E+m$ with $m\ge2$, so it has the step $p_E+1$.
(2) For $2\le i\le m-1$, $H(E_i)=H(L_i)$. So $E_i$ and $L_i$ are taken at the same class-0 rhombus and remove the same tiles. $E_i$ is the $T_E$-configuration at that rhombus whose $G$ equals $G(L_i)=G^{T_L}_t(R)=G^{T_E}_t(R)$ (Lemma~\ref{lem:c0}), with $t$ the a-type of $L_i$. Distinct types have distinct $G$ (by Lemma~ovl the square of one is not a tile of the other). So $E_i$ has type $t$, and by Lemma~\ref{lem:c0} it is the same step as $L_i$. At $i=m$, (F3) says both are $(0,b)$ on the same hexagon. Then $k=m-2$ in Theorem~\ref{thm:corr}, which is defined by the first b-step at a shared rhombus.
(3) By Theorem~bsall, the re-run's middle is $(C-\alpha)\cdot(a+\gamma)$, and the two journeys agree outside the middles. The first steps of the middles, $a=L_0=(1,a_{T_L})$ and $L_1-\alpha$ (a $(1,b)$), differ. The last steps, $L_m$ (a $(0,b)$) and $a+\gamma$ (a $(1,a_{T_L})$), also differ. So the middle in the sense of the G3$'$ tex (between the longest common prefix and suffix) is $a\cdot C$, of length $m+1$.
\end{proof}
\noindent So in case~b the corridor half of (Merge) has \emph{two} proofs. One is (F3) plus Lemma~\ref{lem:c0}. That proof was available to c1: the only missing piece was the one-line finite Lemma~\ref{lem:c0}, and c1 read (Merge) as needing ``target-independence of the engine's choice''. The other is Theorem~\ref{thm:corr}. It is local, and it does not use the sheet-tex chain behind (F3). Each confirms the other. Census: (M1) and (M2) hold 962/962 (\texttt{corr5.out}, \texttt{corr6.out}).

\section{The re-run front, and (H1) in every case}\label{sec:rerun}
\begin{proposition}[$E$'s re-run at the front: presence and Lemma~U]\label{prop:rerun}
Assume (S), including (Mos), so that Lemma~U applies on $M_0\cdot J_L$. Let the contact pair $(E_0,L_0)$ be one of the four of Lemma~\ref{lem:loc2}, and take $L_1$ from Theorem~bsall: $(1,b)$ in case~b, $(1,a_{T_E})$ in case~a. Put
\[K=\bigl(C(L_0)\cup C(L_1)\bigr)\setminus\{\text{rhombi}\}\ \cup\ \bigl(N_E\setminus(N_L\cup G(L_1))\bigr),\]
where $N_E,N_L$ are the non-rhombus tiles of $H_E,H_L$. Then:
\begin{enumerate}
\item $K$ lies on $E^L$'s board $Y_p$ before step $p_E$.
\item Exactly one type $t'$ has $G^{T_E}_{t'}(R_E)\subseteq K$, so $E^L_0$ is that configuration: $(0,b)$ if $E_0=(0,a_{T_E})$, and $(1,a_{T_L})$ if $E_0=(1,b)$.
\item If $E^L$ does not stop at the rhombus $R'$ created by $E^L_0$, then $E^L_1=(1,b)$. It is the unique type whose tiles lie in $K'=(K\setminus G(E^L_0))\cup C(E^L_0)$.
\end{enumerate}
\end{proposition}
\begin{proof}
(1) Created tiles: by Lemma~\ref{lem:fin} no non-rhombus tile created by $J_L$ is removed later by $J_L$. No configuration removes a tile it creates (\texttt{recreate.out}), and $J_L$ never re-creates a tile it removed (Lemma~nr). So these tiles are in $\mathrm{Xadd}_L=M_0\cdot J_L\setminus M_0$. By Lemma~dinv (applied to $J_E$ and $E^L$, identical before $p_E$), $Y_p\setminus X_p=\mathrm{Xadd}_L$, so they lie on $Y_p$.
The $M_0$ tiles: $h\in N_E$ lies on $X_p$, since $E_0$ removes it. By Lemma~dinv, $X_p\setminus Y_p=W_L$, so it suffices that $J_L$ never removes $h$. Go through $J_L$'s steps. Its prefix does not meet $W_E\ni h$ (minimality of $p_L$). $L_0$ and $L_1$ remove $h$ only if $h\in N_L\cup G(L_1)$, which is excluded. The rest of $L$'s middle consists of $L_2,\dots,L_m$ in case~b; in case~a the middle is $L_0L_1$. In case~b these steps remove $H(E_2),\dots,H(E_m)$ (Proposition~\ref{prop:shadow}), and $J_E$ removes $h$ only at $E_0$ (Lemma~\ref{lem:fin}). After the middle, $J_L$'s steps are $L^E$'s (Theorem~bsall). $L^E$ runs on $M_0\cdot J_E$, where $h$ is absent (removed by $J_E$, never re-created). $L^E$ cannot remove a tile it created (Lemma~\ref{lem:fin}). So $L^E$ never removes $h$.
(2) The finite read \texttt{rerun2.py} finds exactly one such $t'$ in each of the eight classes, with the stated values. Its tiles come from $C(L_0)$, $C(L_1)$ and $N_E$, one each, which is c1's trail table. $E^L$ takes a step at $p_E$ (Theorem~fc). By Lemma~U on $Y_p$, the present good configuration is unique and is the one taken.
(3) The same read at $R'$ with $K'$, which lies on $E^L$'s next board by (1) and by the step $E^L_0$. Exactly one type, $b$, is present, with tiles from $C(L_0)$ (two) and $C(L_1)$ (one). Then Lemma~U applies. $R'$ has class~1 in all eight classes.
\end{proof}

\begin{theorem}[(H1) holds]\label{thm:h1}
Assume (S). In every diverging $B/C$ pair, $J_E$ takes a step at $p_E+1$, and $E^L$ takes a step at $p_E+1$. Hence Theorem~\ref{thm:share2} holds with no hypothesis beyond (S), and Proposition~\ref{prop:rerun}(3) applies: $E^L_1=(1,b)$.
\end{theorem}
\begin{proof}
Use the stop grid of the case-a tex (Lemma~grid and Lemma~gridmono there; $F^T=(F^T_0,F^T_1)$). A journey toward $T$ that stops at $R$ has $F^T(R)=f^T(X_T)=:c$. $F^B$ is non-decreasing along $B$-journeys and $F^C$ non-increasing along $C$-journeys, with the increments of the six configurations listed there. $J_E$ (on $M_0$) and $E^L$ (on $M_0\cdot J_L$) both stop for nest $T_E$, which $J_L$ does not touch. So both end at $F=c$, and both start their middles at $R_E$. Suppose one of them, $X$, stops right after its step at $R_E$, while the other, $Y$, has taken the steps $s_1,\dots,s_r$ from $R_E$. Then $c-F(R_E)=\Delta F(X\text{'s step})$, and $c-F(R_E)\ge\sum\Delta F(s_i)$ for $B$ ($\le$ for $C$). \texttt{h1grid.py} evaluates these with the table on all eight classes, in this order:
\begin{itemize}
\item $T_E=C$. First, $E^L$ does not stop at $R'$: the steps $s$ are just $E_0$, and the inequality fails. So Proposition~\ref{prop:rerun}(3) gives $E^L_1=(1,b)$. Then $J_E$ does not stop at $R^E_1$, against the steps $E^L_0,E^L_1$. For example, $E_0=(0,a0)$ gives $(-1,0)$, while $E^L_0E^L_1=(0,b)(1,b)$ gives $(-1,-1)$, and $(-1,0)\not\le(-1,-1)$.
\item $T_E=B$. First, $J_E$ does not stop at $R^E_1$, against $E^L_0$ alone: $(0,1)\not\ge(1,1)$ for $E_0=(0,a1)$, and $(0,0)\not\ge(1,0)$ for $E_0=(1,b)$. This is (H1), so Theorem~\ref{thm:share2} gives $E_1$. Then $E^L$ does not stop at $R'$, against $E_0,E_1$, and in case~b also $E_m=(0,b)$ (Proposition~\ref{prop:shadow}).
\end{itemize}
All eight combinations are excluded (\texttt{h1grid.out}, \texttt{ALL EXCLUDED: True}). Neither ordering uses a fact it is meant to prove. A journey that does not stop takes a step, because it terminates and is never stuck.
\end{proof}
\noindent\emph{First attempt, recorded.} In my first version I used only $E_0,E_1$ for $T_E=B$ in case~b, and the check failed there: $E_1=(0,a0)$ has $\Delta F^B=0$. Adding $E_m$ closed it.

\begin{corollary}[(W$_F$) is proved]\label{cor:wf}
Under (S), in case~b, the front window $(E_0,E_1,E^L_0,E^L_1,L_0,L_1,L^E_0)$ is a translate of one of the four (W$_F$) classes. In case~a, the prefix $(E_0,E_1,E^L_0,E^L_1,L_0,L_1)$ is a translate of one of four classes.
\end{corollary}
\begin{proof}
$E_0,E_1,L_0,L_1$: Theorem~\ref{thm:share2} with Theorem~\ref{thm:h1}. $E^L_0,E^L_1$: Proposition~\ref{prop:rerun} with Theorem~\ref{thm:h1}. $L^E_0=L_1-\alpha$: Theorem~bsall, with $\alpha$ fixed by $L_0$'s type. Each piece is determined by the contact class, so the tuple is.
\end{proof}
\noindent\emph{Census} (\texttt{rerun\_census.out}, \texttt{rerun2.out}, last line). The classes of $(H(E_0),H(E^L_0),H(E^L_1),H(L_1))$ in the shared-square frame, over all 779 + 718 records, are exactly the 8 local predictions, tile for tile. $E^L$ has steps $p_E$ and $p_E+1$ in every record.

\section{What this changes, and the next link}\label{sec:next}
\begin{itemize}
\item \textbf{(E-struct)}, item~2 of the G3$'$ gap (``the shape of $E_1$''), is \textbf{closed} under (H1). The remark there that ``for $E$ nothing does yet'' was wrong in its expectation. No statement about $M_0$ beyond the two contact hexagons was needed. The tiling fixes $E_1$ once $L_1$'s neighbourhood is in the same $M_0$.
\item \textbf{(Merge)}, item~3 of the G3$'$ gap, is \textbf{proved on the corridor} (Theorem~\ref{thm:corr}(1)). c1's ``target-independence of the engine's choice'' turns out not to be a property of the engine's choice at all. The removed sets at a class-0 rhombus are target-blind (Lemma~\ref{lem:c0}). So is the stopping test (Lemma~\ref{lem:pack}). $M_0$ being a tiling then forces the types to agree (Lemma~ovl). No statement about the boards' difference was needed, because both journeys remove only $M_0$ tiles.
\item \textbf{What remains of (Merge).} (M1) (the exit is a step) and (M2) ($k=m_L-3$) are proved in case~b (Proposition~\ref{prop:shadow}). Lemma~red uses only the corridor identity on $2\le i\le m_L-2$, which is now proved. Still computed: (M3), no shared step after the exit (B/C 962/962; it fails 249/249 on $B/B$). (M3) is the target-sensitive clause, and Lemma~red does not use it.
\item \textbf{(H1)} is proved in every case (Theorem~\ref{thm:h1}). The remark below, written before \S\ref{sec:rerun}, is kept as the route that failed. It was: The (Ea) grid argument of the 11-23 c2 tex goes part of the way. Suppose $J_E$ stopped at $R^E_1$. $E^L$ diverges from $J_E$ at $R_E$ with a different good $T_E$-step (Theorem~fc), terminates, and stops on nest $T_E$'s grid. Nest $T_E$ is untouched by $J_L$. Monotonicity of $F^{T_E}$ along $E^L$ (Lemma~gridmono there) then forces $\Delta F^{T_E}(E_0)\ge\Delta F^{T_E}(E^L_0)$ componentwise for $T_E=B$, and $\le$ for $T_E=C$. Using that table:
(i) $T_E=B$, $E_0=(1,b)$, with $\Delta F=(0,0)$: $E^L_0\in\{(1,a0),(1,a1)\}$, with $\Delta F=(1,0)$ or $(0,1)$. Contradiction either way. So \textbf{(H1) is proved in this subcase}, without knowing $E^L_0$.
(ii) $T_E=B$, $E_0=(0,a1)$, with $(0,1)$: this is a contradiction if $E^L_0=(0,b)$, with $(1,1)$, but not if $E^L_0=(0,a0)$, with $(0,0)$. So it needs the open re-run column $E^L_0$, which the census gives as $(0,b)$.
(iii) $T_E=C$: for $E_0=(0,a0)$, with $(-1,0)$, and $E_0=(1,b)$, with $(-1,-1)$, the inequality $\le$ holds against every alternative. So nothing follows.
So (H1) in case~a remains open for $T_E=C$ and for $E_0=(0,a1)$.
\item \textbf{Next unproved links of (E-struct):} (Merge$'$), the corridor on the re-run boards. The boards' initial tiles differ, so the $M_0$-tiling argument does not apply as it stands. Its base case is Corollary~\ref{cor:wf}, and the presence argument of Proposition~\ref{prop:rerun} is the natural tool. Then (W$_B$), and (LEN) for $E$. The re-run columns of (W$_F$) are now proved (\S\ref{sec:rerun}).
\end{itemize}

\section*{Status}
\begin{itemize}
\item \textbf{Proved:} Lemma~\ref{lem:fin} (finality). Corollary~\ref{cor:m0}. Lemma~\ref{lem:loc2} (finite, exhaustive, exact). Theorem~\ref{thm:share2} ((Share$_2$), under (S) and (H1)). Corollary~\ref{cor:front}.
\item \textbf{(H1)} is \textbf{proved} in every case (Theorem~\ref{thm:h1}), so Theorem~\ref{thm:share2} needs only (S).
\item \textbf{Proved:} Proposition~\ref{prop:rerun} (finite, exact), Theorem~\ref{thm:h1}, Corollary~\ref{cor:wf}: (W$_F$) is proved.
\item \textbf{Proved:} Proposition~\ref{prop:shadow}: in case~b, (H1), (M1), (M2), and the corridor identity a second time, from (F3).
\item \textbf{Computed:} the Step~0 tables, size~5 full and size-6 chunks 11--12.
\item \textbf{Negative:} the zone falsifier (\S\ref{sec:step0}(b)).
\item \textbf{Not separating:} (Share$_2$) holds in the local $B/B$ classification and on 440/440 $B/B$ census pairs. The separating clause is (M3).
\item \textbf{Proved (\S\ref{sec:corr}):} Lemma~\ref{lem:c0} (finite), Lemma~\ref{lem:pack} (engine rule (S4)), Theorem~\ref{thm:corr} (the corridor), under (S) and (H1).
\item \textbf{Computed:} (M3) of \S\ref{sec:next}.
\item \textbf{Next unproved link:} (Merge$'$), the corridor on the re-run boards, whose base case is now Corollary~\ref{cor:wf}; then (W$_B$), (LEN) for $E$, and (W$_a$) beyond the prefix.
\end{itemize}
\end{document}
